% =====================================================================
%  R. Nielsen, OM SINDSFORMØRKELSE, BIDRAG TIL SELVMORDERES PSYKOLOGI
%  In: For Ide og Virkelighed. Et Tidsskrift, udgivet af R. Nielsen, Bj. Bjørnson og Rud. Schmidt.  Aargang 1869.
%  Kjøbenhavn: Chr. Steen & Søns Forlag; Louis Kleins Bogtrykkeri.  pp. 689-706.  Antiqua.
%
%  TRANSCRIPTION -- GENERATED; DO NOT EDIT BY HAND.  The five Nielsen essays of the 1869 volume are built by ONE pipeline
%  (ebook-method/for-idee-1869/) into a combined working file, which split5.py cuts into the five transcriptions.
%  Rebuild from the repo root:   sh ebook-method/for-idee-1869/rebuild.sh
%
%  STATUS: COMPLETE (2026-10-13): the volume's title page and the whole essay, pp. 689-706, with its notes.
%
%  EDITION.  The only printing.  Text of record: the Internet Archive's scan of the University of Toronto copy
%  (forideogvirkelig00kjobuoft), ~/bibliotek/Nielsen, Rasmus/1869-for-idee-og-virkelighed.pdf (SCANS.tsv), 886 PDF pages,
%  400 ppi colour; text layer = the IA's ABBYY FineReader OCR.  Printed page N = PDF N + 10 to p. 688; a plate is bound in
%  before p. 689, so pp. 689-706 = PDF N + 12.  Title page PDF 7.  No e-book, no other scan, no earlier transcription found
%  (checked 2026-10-13: the repo, the web, the Internet Archive; HathiTrust and Google Books could not be queried that day).
%
%  ACCURACY (2026-10-13; one collation for all five essays, 189 pp., ~48,000 words).  The words are those of the ABBYY layer
%  (draft.py), collated word by word with an independent tesseract reading of the images (collate/check.py, witness 'ocr'):
%  1132 disagreements at the first collation.  check.py's own rules settled 96; rules.py settled ~330 more, each rule
%  switched on only after a BLIND PILOT of 107 crops (reader A, 21/22 controls) had measured it on this book: tesseract's
%  glyph faults 8/8, aa read an 8/8, accents it adds 8/8, the layer's glyph faults 7/7 and 8/8, o/ø by the ø rule 8/8, the
%  printer's space before a mark 8/8 (house rule), and the narrow word-division rules; the quote, dash, note-number and
%  word-seam rules failed the pilot and stayed off.  The rest were read at the image BLIND (sheets.py: 579 word crops on
%  58 sheets, two planted controls a sheet, 116 of 116 controls right; six readers) and about 110 spots in the main
%  session from line crops (zoom.py).  STRUCTURE read at the image (reader K): the essays' headings, the volume's title
%  page, every footnote whole (the notes were collated where they stand and then moved to their marks), the signatures.
%  LETTERSPACING (the print's emphasis, \emph): candidate lines by box width (spacedwidth.py; the gap measure is noise on
%  this scan): 163 lines read blind (22/22 roman controls), doubtful ones read in the main session.  ITALICS by stroke
%  slant (check.py --italics): 24 crops read blind; the only italics are Latin quotations.
%  SWEEPS after OPEN 0, each spot read at the image: the W-shortens / T-longer list of the blind decisions (found two
%  W decisions that had dropped a note's mark and a closing «, and one that had joined „der f.“); --moved and --splits;
%  the split sweep; the reader-line sweep; paragraph starts against tesseract's indents; quotation marks (the print
%  closes with « and often with », both as printed); the unknown-word sweep (31 spots: 16 corrections).
%  TRANSLATION as a last reading: the twelve translators listed 79 odd spots (ebook-method/for-idee-1869/tr/odd-*.tsv),
%  ~60 read at the image (zoom zo, zw, zx): mended -- f's whose head did not ink, read t by both OCRs (pp. 21, 36, 309, 319,
%  441; a recurring defective sort, as in Folkelige Foredrag), two blind decisions that had taken a reader's t for the
%  print's f (pp. 322, 539), a speck read as a stop (p. 589), a lost dash (p. 525), the verse on p. 414 set as one quotation,
%  false paragraph breaks at page turns left by the notes' removal (pp. 430, 553, 620, 695, 703); the rest are the print's.
%  SAMPLED PAGE CHECK (pagecheck.py, seed 202610136): pp. 12, 418, 521, 694 and the title page read whole and blind, 1069
%  words: 0 word and 0 punctuation errors; one letterspaced passage in a note missed (p. 694), mended.
%  What this cannot see: a misprint both OCRs normalise into a real word.
%  The ø rule (Hans 2026-10-08): o-forms that are no word (Oieblik, Ifolge, imode ...) are set with ø.
%  PRINTED AS IS p. 691: „Tresinstyvende“
%  PRINTED AS IS p. 694: the note mark is printed „¹;«.“ (its parenthesis damaged or a semicolon): set as the note
%  DEFECT, letter supplied p. 695: the note mark „¹)“: its parenthesis damaged
%  PRINTED AS IS p. 695: note: „(af Taillandier.“, the parenthesis not closed; „“Fædrelandet“s“ with a high opening mark
%  PRINTED AS IS p. 697: „færreste. den“: a stop before a lower-case word
%  PRINTED AS IS p. 701: note: „Synsmaade. at“, „sig. kan“, „ramt. Teppet“: stops where the sense wants commas
%  PRINTED AS IS p. 703: note: „Sygdom. til Døden.“
%  PRINTED AS IS p. 704: „mere?»“: the closing mark printed »
%  Not transcribed: running heads, folios, sheet signatures („For Ide og Virkelighed.  20“), the library's stamp; the large
%  initial of each essay is set as a plain letter.  Quotation marks as printed (» « in the text, „ “ in the notes).  The notes
%  are numbered ¹) per page, as printed (p. 553: *)).
%
%  TOOLS (ebook-method/for-idee-1869/, copied from folkelige-foredrag): draft.py, ocr.py, rules.py, sheets.py + reads.py +
%  READER-BRIEF.md, heads.py, spacedwidth.py + spacedlines.py + spaced.tsv, zoom.py, pagecheck.py + PAGECHECK-BRIEF.md,
%  patches.py (PATCHES, NOTES, POST), decisions.tsv, the sweeps, split5.py, rebuild.sh.
% =====================================================================
\documentclass[12pt,a4paper,openany]{book}
\usepackage[utf8]{inputenc}
\usepackage[T1]{fontenc}
\usepackage{libertinus}
\usepackage{libertinust1math}
\usepackage[danish]{babel}
\usepackage[protrusion=true,expansion=true]{microtype}
\usepackage{setspace}
\usepackage[margin=1.3in]{geometry}
\usepackage{fancyhdr}
\usepackage{hyperref}
\hypersetup{pdftitle={Om Sindsformørkelse}, pdfauthor={Rasmus Nielsen}, hidelinks}
\pagestyle{fancy}
\fancyhf{}
\fancyhead[LE,RO]{\thepage}
\fancyhead[C]{\textit{R. Nielsen: Om Sindsformørkelse (1869)}}
\setstretch{1.3}
\usepackage{marginnote}
\renewcommand{\thefootnote}{\arabic{footnote})}   % the print numbers each page's note ¹) (p. 553: *))
\newcommand{\opage}[1]{\marginnote{\footnotesize\textit{[#1]}}}
% the essay's heading: title (and subtitle), centred, a short rule; the print's large initial is not reproduced
\newcommand{\essayhead}[2]{\begin{center}{\large #1}\ifx&#2&\else\\[0.3em]{\normalsize #2}\fi\\[0.5em]\rule{3em}{0.4pt}\end{center}\medskip\noindent}
\newcommand{\signature}[1]{\par\smallskip\hfill #1\hspace*{2em}\par}

\begin{document}

% --- front matter: the volume's title page (PDF 7), read at the image (reader K, F7; the page check reads it again blind).
\begin{titlepage}
\begin{center}
\vspace*{3em}
{\Huge For Ide og Virkelighed.}\\[1em]
\rule{8em}{0.4pt}\\[1em]
{\large Et Tidsskrift,}\\[0.8em]
{\small udgivet}\\[0.4em]
{\footnotesize af}\\[0.6em]
{\large R. Nielsen,}\\[0.4em]
{\large Bj. Bjørnson og Rud. Schmidt.}\\[1em]
\rule{8em}{0.4pt}\\[1em]
{\large Aargang 1869.}\\[1em]
\rule{8em}{0.4pt}\\[4em]
Kjøbenhavn.\\[0.4em]
Chr. Steen \& Søns Forlag.\\[1.5em]
{\small Louis Kleins Bogtrykkeri.}
\end{center}
\end{titlepage}

% ===== TEXT (pp. 689-706) =====
\essayhead{Om Sindsformørkelse,}{Bidrag til Selvmorderes Psykologi.}
% --- p. 689 ---
\setcounter{footnote}{0}%
\opage{689}Vil Nogen i al Almindelighed opkaste et Spørgsmaal
som dette: hvem er vel den største Forbryder, enten
Morderen slet og ret eller Selvmorderen, forudsat, at begge
have handlet med fuldt Overlæg? da maa det forekomme
den sunde Sands, som kunde Svaret ikke være tvivlsomt.
Der er i Morderens Sind en saa brutal Egoisme,
en saa grusom Ringeagt for Næstens Liv, en saa oprørende
Trods mod guddommelige og menneskelige Love,
at Samfundet i denne blodige Misdæder med Rette seer
sin svorne Fjende. Anderledes med Selvmorderen. En
saadan Ulykkelig blev jo ikke farlig for Samfundet.
Han trængte sig ikke frem med Vold; nei, han trak
sig i al Stilhed tilbage. Han greb ikke Dolken for at
rydde Andre afveien og skaffe sig Plads, han vilde tvertimod
som En, der var tilovers i Livet, selv gaae afveien,
fjerne sig i al Ubemærkethed, for at der kunde blive Plads
for Andre. Vel sandt, han sønderrev formastelig sin
egen Livstraad, men  hvo veed? --- en giftig Braad
af Dødens Orm sad vel allerede dybt i Hjertet. Forestillingen
om, hvad et Menneske maa lide og stride,
inden det i Livets Skibbrud slipper den sidste Planke,
det sidste Haab, og vælger Fortvivlelsens sidste Udvei,
kan hos Betragteren vække alvorlige Tanker.

% --- p. 690 ---
\setcounter{footnote}{0}%
\opage{690}For i hvert paagjældende Tilfælde at komme efter,
hvad der hos den enkelte Selvmorder vel kan have
været Bevæggrunden, er man naturligvis opmærksom
paa de nærmere Omstændigheder og da især paa de
dels tilfældige Yttringer, dels udtrykkelige Erklæringer,
som den Afdøde maatte have efterladt sig. Men det er
netop det Mærkelige, at slige Udtalelser ikke sjelden
staae i et saa skærende Misforhold til det, som deraf
skulde forklares, at de umiskjendelig vise hen til noget
Andet, noget Skjult, der maaske ubevidst eller dog kun
halvbevidst har rørt sig i Sindet. Hvad dette, Gjerningens
egenlige og sande Grund, vel kan have været, det er
en mørk Gaade, en Hemmelighed, som den Døde tog
med sig i Graven. Paa disse Vilkaar faaer da Betragteren
ikke ligefrem med den Hedengangne, men med sig
selv at gjøre. Det Virkelige opløser sig for ham i lutter
Muligheder; jo længere han stirrer ned i disse Muligheders
Afgrund, desto mere udsætter han sig for at
svimle; en Tanke stiger op fra Dybet og bringer Blodet
til at isne, en mørk Tanke med det Spørgsmaal:
kunde Sligt ikke ogsaa hændes Dig? Jo mere Idealitet,
der er i Betragterens Væsen, desto naturligere vil det
falde ham at udbrede et formildende Skær endog over
det raaeste Selvmord.

Vi prise det Menneskelige i en saadan Medfølelse;
men derfor tør vi dog ikke glemme, hvad et Selvmord
er i sig selv, at det er en \emph{Forbrydelse}. Vi tør saameget
mindre glemme det, som Hang til denne Forbrydelse
endog i betænkelig Grad er udbredt hos os.

Hvor hyppigt Selvmord gjennemsnitlig er i Danmark,
kan skjønnes af »den Oversigt over Forholdet af
% --- p. 691 ---
\setcounter{footnote}{0}%
\opage{691}Selvmord til Dødsfald i Almindelighed i Aldersklasserne
over 10 Aar, 1835 --- 1845,« som Dr. med. Keyser (Dansk
Folkeblad 1846) har meddelt. Forfatteren tilføier: »Man
seer heraf, at i hele Danmark er hver Tresinstyvende   % PRINTED AS IS (p. 691): „Tresinstyvende“
af de Døde over ti Aar en Selvmorder, eller med andre
Ord, at næsten to Procent af alle Dødsfald ere foraarsagede
herved; at denne Hyppighed stiger i Kjøbenhavn;
at den for Mandfolk i denne By naaer en foruroligende
Høide, idet hvert syv og tyvende døde Mandfolk
over ti Aar er en Selvmorder, og at altsaa i denne
Alder næsten fire Procent af de Afdøde bortrykkes paa
denne Maade.« Denne Kjendsgjerning maa i høi Grad
forbause Enhver, der ikke selv har havt Leilighed til
at undersøge Sagen; den maa, hvad Forfatteren da ogsaa
fremhæver, minde om den store Vanskelighed, der
er forbundet med at gjøre en paalidelig Erfaring, og
»hvormeget Uvæsen der bliver drevet med dette Ord».

I »\emph{Stadier paa Livets Vei}«, et Skrift, der udkom
1845, udtaler Pseudonymen sin Glæde over at
være Forfatter i et lille Land, (Side 378), og siger:
»Nogle af mine Landsmænd mene vel, at Kjøbenhavn
er en kjedelig By og en lille By. Mig synes den tvertimod,
forfrisket af det Hav, ved hvilket den ligger, og
uden selv i Vinteren at kunne opgive Bøgeskovens Erindring,
at være det lykkeligste Opholdssted, som jeg
kjender. Stor nok til at være en større By, lille nok
til, at der ingen Torvepris er paa Mennesker; kan den
tabellariske Trøst, som man har den i Paris over saa
og saa mange Selvmordere\dots{}. ikke trænge sig forstyrrende
ind« osv. Skjønt vi nu villig indrømme,
at Digterstemninger og statistiske Tabeller ere inkom%
% --- p. 692 ---
\setcounter{footnote}{0}%
\opage{692}mensurable Størrelser, saa holde vi os dog overbeviste
om, at Hentydningen til den tabellariske Trøst over saa
og saa mange Selvmordere vilde være falden anderledes
ud, hvis »Stadiernes» Forfatter havde havt en Forestilling
om det Resultat, hvortil en »sammenlignende
Oversigt over Hyppigheden af \emph{Selvmord} i Danmark,
1835—44, og i Frankrig, 1835—43«, virkelig har ført;
thi Resultatet er, at \emph{Selvmord omtrent er tre
Gange saa hyppigt i Danmark som i Frankrig,}
og at \emph{Hyppigheden} af Selvmord saavel for hvert
enkelt Kjøn som for dem begge tilsammen, er \emph{større}
i \emph{Kjøbenhavn} end i hele \emph{Seinedepartementet}.

Den, der forsker efter Selvmordets almenpsykologiske
Grundbetingelse, støder overalt paa store Vanskeligheder.
Vistnok kan Forskjel i Alder og Kjøn udøve
Indflydelse, men hvilken? Det Udgangspunkt, der
søges i Temperamentsforskjellighederne, er høist utilforladeligt,
thi Erfaringen lærer, at Sangvinikeren og
Flegmatikeren ligesaa vel som Kolerikeren og Melankolikeren
kunne findes paa Selvmordernes Liste. Ingen
Forskjel mellem Syge og Sunde, mellem Feige og Modige,
mellem Dannede og Udannede, mellem Rige og
Fattige, ja ei engang mellem overvættes Lykkelige og
grændseløs Ulykkelige kan i dette gaadefulde Punkt
opstilles med tilbørlig Sikkerhed.

At en Mand med trodsigt Mod kaster Livet bort,
naar det bliver ham en Byrde, finder man naturligt,
men at Feighed og Angst kunne føre til samme Resultat,
maa synes besynderligt. Dog haves Exempler. Ligesom
det for Mange er en livsfarlig Sag at sidde selvtrettende
til Bords, saaledes ansees det ogsaa af Almuen
% --- p. 693 ---
\setcounter{footnote}{0}%
\opage{693}paa Landet for livsfarligt, at En under Jordpaakastelsen
falder ned i Graven; thi den, der falder ned til den
Døde, maa døe inden Aaret er omme. En Mand, hvem
Uheldet ramte, tog sig det saa nær, at han hængte sig
et Par Maaneder efter: af Frygt for Døden valgte han
at døe.

Et Selvmord kan begaaes i Raseri; men at det
ogsaa kan udføres i al Besindighed, med klar Overskuelse
over Situationen, derom vidner følgende Brev,
hvormed en fredelig Leier har villet berolige sin Vertinde:
»Gode Madam N.! I den Tilstand, De vil finde mig i,
vil De strax raabe og skrige: han har hængt sig, men
det maa De ikke. De behøver jo kun at sige, jeg er
død\dots{} Jeg lader mig begrave fra Hospitalet. Hr. H.
skal have nogle Penge, sig ham, de blive rigtig betalte
af en af mine Venner, som jeg har laant 40 Rdl.
(Folkebl. 1846 S. 72). Denne flegmatiske Uanfægtethed
staaer i skærende Modsætning til den Forbittrelse,
hvormed Haandværksdrengen, (anf. Blad S. 75), som
havde faaet et Ørefigen, tog Afsked fra Verden og skrev
til sine Forældre: »For sidste Gang vil jeg sige Dem
Farvel her i Livet; jeg kan ikke udholde det længere.
Forband Madsen, forband Svenden; to Gange Prygl i
een Nat».

Hvorledes Dannelse og Raahed kunne, trods Bevæggrundenes
Forskjellighed, mødes i den samme Forbrydelse,
er en bekjendt Sag. I anden Del af \emph{Enten
--- Eller} (Anden Udgave S. 172) udkaster Ethikeren
et fra sit Standpunkt træffende psykologisk Livsbillede
af et ungt Menneske, hvem intellektuel Forskruethed
og grublende Indesluttethed drev til Selv%
% --- p. 694 ---
\setcounter{footnote}{0}%
\opage{694}mord. »Skulde Nogen«, skriver den Afdøde i sit Brev til
Justitsraaden, »spørge Dem om Anledningen, kan De sige,
at det var i Anledning af den store Ildebrand. Skulde
Nogen spørge Dem om Tiden, kan De sige, at det var
i den for mig saa mærkelige Juli Maaned\dots{} Ikke anseer
jeg et Selvmord for noget Roesværdigt. Det er
ikke af Forfængelighed jeg har besluttet mig dertil.
Derimod troer jeg paa Rigtigheden af den Sætning, at
intet Menneske kan taale at see det Uendelige. Det
har viist sig engang for mig i intellektuel Henseende,
og Udtrykket derfor er Uvidenhed. Uvidenhed er nemlig
det negative Udtryk for den uendelige Viden. Et Selvmord
er det negative Udtryk for den uendelige Frihed.
Det er en Form af den uendelige Frihed, men den negative
Form. Held den, som finder den positive\footnote[1]{Maaske er Historien digtet. Men at Forskjellen mellem Digt og Virkelighed Intet betyder, hvor Talen kun er om \emph{psykologisk Sandhed}, er en Bemærkning, som den, \emph{der} skriver for dannede Læsere, vel tør ansee for overflødig.}«. Til   % PRINTED AS IS (p. 694): the note mark is printed „¹;«.“ (its parenthesis damaged or a semicolon): set as the note
denne Frihedsstræben i en dialektisk raffineret Dannelses
Form finder man (Folkebl. S. 74) et Modstykke i den
udialektiske Raaheds Form. Det er en Kleinsmedlærling,
der higede mod Friheden i Amerika, men maatte
i Mangel af Pas søge den i en af Frederiksberg Haves
Kanaler, ved hvis Bred, han efterlod Følgende: »Her
ligger Undertegnede lige ud for. Da jeg er kommen
fra min Mester, fordi jeg ikke vilde krybe for Madammen,
og nu da jeg agtede at reise til Amerika, vilde
han ikke skrive paa min Seddel, som han burde, altsaa
kunde jeg ingen Pas faae. Hvad skulde jeg nu gjøre.
Jeg vidste kun dette Raad \dots{}»

% --- p. 695 ---
\setcounter{footnote}{0}%
\opage{695}I hvilken Grad exalteret Geniforgudelse, fantastisk
Geniusdyrkelse, platonisk, hypertrofisk Ideforelskelse kan
angribe Fruentimmernerver, derpaa afgiver den ved sin
Aand og Ynde iøvrigt beundringsværdige Charlotte Stieglitz
et tragisk Exempel\footnote[1]{Jvfr. „Et Drama i det literære Liv“ (af Taillandier. Fortællinger og Skildringer fra “Fædrelandet“s Feuilleton. Oversatte af Rudolf Schmidt. S. 339 flgd.}. Hun ansaae sin Mand for at være   % PRINTER'S DEFECT (p. 695): the note mark „¹)“: its parenthesis damaged   % PRINTED AS IS (p. 695): note: „(af Taillandier.“, the parenthesis not closed; „“Fædrelandet“s“ with a high opening mark
den Digter-Messias, af hvis glandsfulde Optræden hele
Tydskland --- altsaa hele Verden --- maatte vente sin poetiske
Forløsning; det var hende derfor ikke nok at være
hans Hustru, hun maatte ogsaa være hans Muse. Forventningerne
skuffedes, Geniet svigtede, Manden sygnede;
den til Muse udkaarede Hustru satte nu sit eneste
Haab til, at en stærk Sindsrystelse maaske kunde gjengive
den syge Mand sin tabte Spændkraft og vække
hans slumrende Genius; for at fremkalde en saadan aandelig
Rystelse og vie sin Digter ved Lidelsens Daab,
smykkede hun sig som Dødens Brud og stødte saa Dolken
i sit Hjerte.

Den Sjælerystelse, hvorved Charlotte Stieglitz vilde
opløfte sin idesyge Ægtefælle til Idealets Herlighed, staaer
i skærende Kontrast til den Sønderknuselse, som en hevngjerrig
Tjenestepige ved sin Opoffrelse øiensynlig har
tiltænkt sin troløse Elsker. »Hr. Sivertsen! Nu kan De
see, hvad De har for Deres slette Handlinger imod mig,
som for 15 Uger siden har lidt saameget for Deres
Skyld. Og De vil vise Dem saa fræk at sige, at De
ikke er Fader til mit Barn, som Gud i min Dødsstund
skal være Vidne til; men De er et slet Menneske, som
% --- p. 696 ---
\setcounter{footnote}{0}%
\opage{696}det aldrig kan gaae godt, for De udøver saameget Ondt
imod Deres Medmennesker. De gav mig gode Løfter,
men jeg saae ikke Meget af det. Farvel, slette Menneske«.
(Folkebl. S. 75).

Det kunde ved et flygtigt Øiekast næsten synes
upassende at sammenstille to saa ulige Væsener, som
hin ideberusede Sværmerinde med denne vistnok ingenlunde
aandrige Tjenestepige; og dog, naar Sagen skal
frem for Psykologiens og Ethikens Domstol, er det et
stort Spørgsmaal, om Idevildfarelsen ikke i Evighedens
Forstand har langt farligere Følger end den raa Forvildelse.
Naar Bevæggrundene ere saa forskjellige, maae
Handlingerne vel ogsaa være det. Men Handlingernes
Ulighed taber sig her i Gjerningens Lighed: det er i
Virkeligheden samme Gjerning, samme Forbrydelse.

De almindeligste Bevæggrunde til Selvmord har
Dr. Keyser (Folkebl. S. 52) henført under følgende syv
Kategorier: Næringssorger, ulykkelig Kjærlighed, Familiesorger,
legemlig Sygdom, Misfornøielse med sin Stilling
og med Livet, slet Opførsel og Sindssvaghed. Forsaavidt
denne Klassifikation retter sig efter faktiske,
navnlig franske, Angivelser, kan den vel i tabellarisk
Henseende have sin Gyldighed; men fra et psykologisk
Standpunkt er Kategorien: Misfornøielse med Livet,
aabenbart for almindelig, og Kategorien: Sindssvaghed,
naar Ordet tages i sædvanlig Betydning, saa speciel, at
den udelukker Tilregnelighed. Mærkelige blive de anførte
Bevæggrunde fornemmelig ved den Maade, hvorpaa
de forbinde sig med Aarstidernes Indflydelse.

At Aarstidernes Vexel, ligesom Afvexling af Dag
og Nat, indvirker paa Sindets Stemning, er en be%
% --- p. 697 ---
\setcounter{footnote}{0}%
\opage{697}kjendt Sag. Man skulde nu mene, at den lyse og
varme, paa ydre Livsbetingelser rigeste Del af Aaret
maatte tælle de færreste. den mørke, kolde, af Næringssorger   % PRINTED AS IS (p. 697): „færreste. den“: a stop before a lower-case word
tyngede Aarstid de fleste Selvmord; men,
mærkværdigt nok, Forholdet er just det omvendte:
December og Januar have de færreste, Juni og Juli de
fleste. Ligesaa gaadefuld er Aarstidernes Indflydelse
paa Dødsmaaden. Vi kunne til Nød forstaae, at Hængningen
i det Hele foretrækkes om Vinteren, da Vandet
er saa koldt, og at Drukningen har noget Tillokkende
i Sommerheden; men det Uforklarlige ligger i Overgangen.
»Efter at Hængning i December Maaned er
sunket ned til sit Laveste, stiger den pludselig til sit
Høieste i Januar, og efterat Drukning har naaet sit
Høieste i December, synker den pludselig ned til sit
Laveste i Januar.« Dersom dette er mere end et Spil
af Tilfældet, da maae Aarsagerne være saa skjulte, at
Ingen endnu kan ane, hvori de bestaae; kun Saameget
kan siges, at ingen Regel danner sig uden en Kreds af
staaende Betingelser, og at ingen Virkning er uden
Aarsag.

Søge vi et Enhedsudtryk for den sjælelige Tilstand,
hvori Selvmorderen i Gjerningens Øieblik maa antages
at befinde sig, da er der vel neppe noget enten mere
omfattende eller mere betegnende Ord end \emph{Sindsformørkelse.} Ligesom Nattens Mørke gjør Øiet blindt
og lader Alt løbe i Et, saaledes kan ogsaa Mørket
i Sjælen lade de allerforskjelligste Bevæggrunde løbe
sammen; Fornuftens Lys kan da paa en Maade siges
at være gaaet ud, før det fysiske Liv udslukkes. Det
Punkt, hvortil den psykologiske Interesse fæster sig,
% --- p. 698 ---
\setcounter{footnote}{0}%
\opage{698}er Sindsformørkelsens Oprindelse, dens langsomme Stigen,
dens pludselige Indtræden.

Hvorfra kommer Sindsformørkelse? See vi bort fra
den Svaghed, der nærmest maa tilskrives organiske Forstyrrelser i Hjernen, saa tør det vel antages, at den kommer
indenfra, fra Sindets Grund, fra et hemmeligt Mørke,
som har unddraget sig Bevidstheden. Ligesom hele Jordkloden
aldrig er oplyst paa een Gang, men den ene
Halvdel altid har Nat, medens den anden har Dag, saaledes
er Sjælens Liv ogsaa delt mellem Lys og Mørke.
Dog, Lignelsen halter; thi medens Kloden ved at dreie
sig om sin Axe successivt bliver belyst paa alle Sider,
er der en Side af Sjælens Liv, den organiske nemlig,
hvor det altid er Nat; og medens Jordens Indre for
bestandig udelukker Sollyset, er Sjælens Indre netop
bestemt til at aabne sig for Bevidsthedens Lys. Sjælens Indre
er Selvet, og Selvbevidsthed bestaaer just deri, at
Mennesket klarer sig sit oprindelige Selv. Det Mørke,
der danner Bevidsthedslivets naturlige Grændse, hører
med til vor Organisation og indjager ingen Frygt; men
Mørket i Selvet, det aandelige Mørke, der unddrager
sig Bevidstheden og dog er bestemt til at vige for
Bevidstheden, er som et Underverdenens Forbud, den
Livets Nat, der huser de farlige Muligheder.

Det aandelige Mørke kan ruge i Sindet, og endda
kan det synes, som var det høi lys Dag i Bevidstheden.
Tingenes Verden med dens Fænomener, Tilværelsen
med dens Opgaver og endelige Interesser, hele
den ydre Virkelighed kan staae for Anskuelsen i fuld
Klarhed, Forstanden kan være i høi Grad udviklet og i
mange Retninger vel orienteret, og dog er Mørket der%
% --- p. 699 ---
\setcounter{footnote}{0}%
\opage{699}inde. Saa kan et saadant Menneske leve med Andre,
arbeide og nyde som Andre, frygte og haabe, lee med
de Glade og græde med de Bedrøvede: der er Ingen i
Omgivelsen, der mærker Noget; ja den Paagjældende
selv har knap en Forestilling om, hvorledes hans Tilstand
er, at han sjælelig talt gaaer med Døden i Hjertet.
Nu og da har han maaske fornummet en uvilkaarlig
Indskydelse; det var et Rædselssyn, en Fantasiens
Lygtemand, en vildsom Tanke, der kom og svandt; men
det var ogsaa det Hele. Dette er Dispositionen.

Om denne Disposition gjælder, hvad S. Kierkegaard
i sit Skrift: \emph{Sygdommen til Døden} S. 17 anfører om
Fortvivlelsen: »Som vel Lægen maa sige, at der maaske
ikke lever et eneste Menneske, der er ganske sund,
saaledes maatte man, hvis man ret kjendte Mennesket,
sige, at der ikke lever et eneste Menneske, uden at han
jo dog er lidt fortvivlet, uden at der dog inderst inde
boer en Uro, en Ufred, en Disharmoni, en Angst for et
ubekjendt Noget, eller for et Noget, han end ikke tør stifte
Bekjendtskab med,\dots{} en Angst for sig selv, saa han dog,
som Lægen taler om at gaae med en Sygdom i Kroppen,
gaaer med en Sygdom, gaaer og bærer paa en Aandens
Sygdom, der en enkelt Gang glimtvis, ved og med
en ham selv uforklarlig Angst, lader sig mærke med, at
den er derinde«. Dispositionen er almindelig, saa almindelig,
at den endog kan antages at være naturlig for det
naturlige Menneske; men i hvilken Grad den er tilstede
hos de Enkelte, om det sjælelige Mørke skal svinde for Aandens
Lys, eller det som en Nattens Hemmelighed engang
i et ulyksaligt Øieblik skal bryde frem af sit Skjul og formørke
det hele Sind, vil afhænge af, hvorledes Menne%
% --- p. 700 ---
\setcounter{footnote}{0}%
\opage{700}sket omgaaes med sig selv, og derhos beroe saavel paa
ydre Omstændigheder som paa indre Betingelser.

For at trænge det aandelige Mørke ud af Sjælen
kræves to Betingelser: væsenlig Selvforstaaelse og villiefast
Karakterdannelse. Hvor en af disse Betingelser
mangler, vil Dispositionen under givne Omstændigheder
let kunne blive saa stærk, at det ender med
Selvmord.

At de Betingelser, paa hvilke den naturlige og dog
saa unaturlige Disposition alene kunne betvinges, kun altfor
ofte savnes, har sin inderste og væsenlige Grund i
Selvforsømmelse, denne igjen i Selvforglemmelse. Selvforglemmelsen
har sine Grader; den kan dels være umiddelbar
og naiv, dels moralsk, dels i strengeste Forstand
religiøs.

Den naive Selvforglemmelse er karakteriseret derved,
at Individet simpelthen opfatter sit Liv som en Naturgave
uden at betænke, at denne Naturgave tillige er en aandelig
Opgave. Den, der i Sandhed skal lære, hvad Livet
er, maa ikke blot have Noget at \emph{leve} af, han maa
i ubetinget Forstand ogsaa have Noget at \emph{leve} for;
det Sidste er ligesaa nødvendigt som det Første. For
den naturlige Sands stiller Sagen sig imidlertid anderledes.
»Havde vi bare nok at \emph{leve} af, saa kunde vi
sagtens faae Noget at \emph{leve} for; thi hvorfor \emph{leve} vi?
Vi skulle arbeide, virke og gjøre Gavn! det er saa Livets
Lov, dets strenge Vilkaar; men det er ikke Livets
Øiemed. Øiemedet er at glæde sig ved Livet, nyde Livet.
Kun den veed at leve, der forstaaer at nyde: jo
større Nydelser, jo flere Fornøielser, desto skjønnere,
rigere og lykkeligere er Livet«. At dette er en
% --- p. 701 ---
\setcounter{footnote}{0}%
\opage{701}meget udbredt og populær Anskuelse, vil vist Ingen
negte\footnote[1]{Hvor naturlig den Synsmaade. at Livets Maal er at nyde og more sig. kan groe fast som Dogme, og hvor uvilkaarlig dette Dogme kan komme til Udtalelse endogsaa hos Ældre, derpaa skulle vi blot anføre et lille Exempel. Et Par adstadige Borgermænd sad en Aften sammen i Theatret og saae Don Juan. Stykket er ude; Lynstraalen har ramt. Teppet falder. Da siger den Ene: „Det var dog en grusom Ende, det tog med ham“. „Ja“, svarede den Anden, „men han havde dog moret sig“.}.   % PRINTED AS IS (p. 701): note: „Synsmaade. at“, „sig. kan“, „ramt. Teppet“: stops where the sense wants commas

Vistnok faaer man engang imellem ligesom en Anelse
om, at Livet muligvis kunde have et andet og høiere
Maal; men paa Selvforglemmelsens første Naturtrin siger
det ikke stort. Ja, kom det an paa Talemaader, kunde
man jo sagtens indvinde en Del høitravende Betragtninger;
men Livet er en Realitet, og i Realiteten hvad er
saa Livet? »Et Pust i Sivet«, en mørk Gaade. Da der
nu ikke kommer Noget ud af at sætte sig hen og gruble
over en saadan Gaade, saa tager man Verden som den
er, nyder saalænge man kan, og morer sig.

Naar saa Livet tynger, naar Udsigterne mørknes,
naar den paa Lykken Stolende, efter Nydelser Higende
rammes af Uheld og trykkes af Modgang, hvad Under
da, at »Næringssorger«, »ulykkelig Kjærlighed«, »Familiesorger«
o. desl., i et skæbnesvangert Øieblik kan
overvælde den Uforberedte, jage Mørket op fra den
hemmelighedsfulde Nat og drive ham til Fortvivlelse!
Jo mere den Anskuelse, at Livet er til for Nydelsens
Skyld, breder sig og fæster Rod i det menige Folk, desmindre
tør vi undre os over, at Selvmord er saa hyppigt,
vi maa snarere undre os over, at det ikke er endnu
hyppigere.

% --- p. 702 ---
\setcounter{footnote}{0}%
\opage{702}Jo mere Nydelsesraseriet breder sig og antænder
Masserne, desto mere voxer den aandelige Umyndighed.
Paa hvilket Umyndighedstrin den naive Selvforglemmelse
kan holde Individerne fangne, er iøinefaldende,
naar vi betragte den Indflydelse, som Aarstiderne udøve
paa den for Liv og Død afgjørende Handling. Det
gaaer med Selvmorderne som med Trækfuglene: det
er ikke en til Frihed modnet Villie, men en dunkel
Naturmagt, der bevæger dem til at tage bort. At det
skeer i omvendt Orden, at Trækfuglene gjerne blive her
om Sommeren og først drage bort, naar Vinteren nærmer
sig, medens Selvmorderne i Reglen foretrække at
overvintre i dette Liv og kun i større Mængde drage
herfra om Sommeren, giver en svag Antydning af, at
Mennesket er Aand. Det er nemlig en psykologisk
Kjendsgjerning, at Sommeren er farligere for Aandens
umyndige Børn end Vinteren. Den haarde Nødvendighed,
den øieblikkelige Trang ægger Selvopholdelsesdriften, opfordrer
til Anstrængelse og forjager Griller; hvorimod
den mildere Aarstid, den smilende Natur med sine dulmende,
bedaarende Indtryk, særlig frister den Ørkesløse
til Sløvhed, Drømmerier og forvirrende Tanker.

For at afbryde Selvforglemmelsen og vække de
Tankeløse, gjør Tilværelsen hvad den kan; som Regn
vexler med Solskin, Vindstille med Storm, saaledes lader
den Held omskiftes til Uheld, Medgang til Modgang, Glæde
til Sorg; den kommer Mennesket imøde, taler igjennem
Tilskikkelserne, og siger til den Forvirrede: gaa i Dig
selv! De, som nu ved slige Paamindelser lade sig vække,
faae forsaavidt en høiere Forestilling om Livets Alvor,
som de erkjende, at det er Villie og Karakter, hvorpaa
% --- p. 703 ---
\setcounter{footnote}{0}%
\opage{703}det kommer an. Livet er en Skole for Karakteren.
Paa dette Trin kan der udfoldes megen Dygtighed, og
Udholdenhed, megen borgerlig og politisk Virksomhed.
Her er Noget at leve for; her er fornuftige Formaal. Et
kraftigt Selv begynder at danne sig. Det er, som man
seer, et første Tilløb til det Ethiske; Karakteren staaer,
saa at sige, i Ethikens Forgaard, men dog, trods sin tilsyneladende
Styrke og Holdning, paa svage Fødder, hvis Sindet
ensidig er udadvendt, udelukkende optaget af Livet
i det Udvortes. Det hele Spørgsmaal om Selvet i dybere
Forstand bliver da, som S. Kierkegaard udtrykker det, »en
Slags blind Dør« i Baggrunden af Menneskets Sjæl, indenfor
den er der Ingenting. Individet »overtager, hvad han
i sit Sprog kalder sit Selv, det vil sige, hvad Evner, hvad
Talenter o. s. v., der kan være givet ham, alt dette
overtager han, dog med Retningen udefter, mod, som det
hedder, Livet, det virkelige, det virksomme Liv; han
omgaaes meget forsigtigt den Smule Reflexion, han har
i sig, han frygter det skal komme op igjen, Det i Baggrunden.
Saa lykkes det ham lidt efter lidt at glemme
det; i Aarenes Løb finder han det saa næsten latterligt,
især naar han er i godt Selskab med andre duelige og
virksomme Mænd, der have Sands og Dygtighed for det
virkelige Liv. Charmant! Han er nu, som der staaer
i Romanen, allerede i flere Aar lykkelig gift, en virksom
og foretagende Mand, Fader og Borger, maaske
endog en stor Mand; hjemme i hans Hus kalde Tjenestefolkene
ham »Han Selv«\footnote[1]{Sygdom. til Døden. S. 53---54.} --- og dog er der i denne   % PRINTED AS IS (p. 703): note: „Sygdom. til Døden.“
»Han Selv«, naar Ethiken foretager sin Synsforretning,
% --- p. 704 ---
\setcounter{footnote}{0}%
\opage{704}en bundløs Dunkelhed, et selvløst Mørke: det evige Lys
mangler. Paa den mørke, selvløse Grund kan en selvisk
fast Villie ikke destomindre reise en Karakterbygning,
der i visse Maader ligner de ægyptiske Templer:
Opgangen er storartet, Murværket imponerende, Søilehallerne
prægtige, men det Allerhelligste er et lille mørkt
Hul. Bærer nu en saadan Individualitet den ulykkelige
Disposition i sig, saa behøves der kun et Slag af Skæbnen,
en Krænkelse, han ikke kan taale, et Pengetab,
han ikke kan bære, et overilet Skridt, der ikke kan
gjøres om, en Legemssvaghed, han ikke vil finde sig i,
og hvad har han nu at leve for? »Det stærke Selv
er sin egen Herre; en Strikke om Halsen, en Kugle
for Panden, et Spring i Søen, hvad er det saa mere?»   % PRINTED AS IS (p. 704): „mere?»“: the closing mark printed »
Saa skeer det. Omgivelsen, der Intet aner, hører pludselig
et uventet Budskab, og spørger overrasket: hvem
skulde have tænkt det, at et saadant Liv skulde faae en
saadan Ende?

Den ethiske Selvforsømmelse har sin dybere Grund
i den religiøse. Vi have i foregaaende Afhandlinger
søgt at paavise den principielle Sammenhæng mellem
Religion og Ethik. Det er en bedaarende Indbildning,
at den menneskelige Villie uden levende Indtryk
af Guds Villie skulde kunne gjøre Ideen Fyldest.

Jo mere Idealitet, der er i et Menneske, desto lettere
udsættes det for Aandsforvildelser, naar den religiøse
Selvforstaaelse mangler. Det er i denne Henseende
ret Betegnende, at Heltinden i hint uhyggelige »Drama«
allerede havde forvildet sig bort fra det Religiøse, før
hun foer vild i Ideerne. Som ganske ung Pige havde
Charlotte Willhoeft, under Indflydelse af en Lærer, som
% --- p. 705 ---
\setcounter{footnote}{0}%
\opage{705}tilhørte Methodisternes Sekt, hengivet sig til en fanatisk
og mørk Fromhed; hun foragtede denne Verden, hun
lod haant om dens virksomme Liv, og mere end een
Gang foer i hendes Higen mod Gud Tanken om Selvmord
gjennem hendes Hjerne. Denne glødende Attraa
efter Døden, denne utaalmodige Higen mod et andet
Liv, var imidlertid kun en Krise; hun vendte tilbage
til Jorden og »bevarede fra sit Samliv med Gud kun
en Kjærlighed til Alt, hvad der er guddommeligt blandt
Menneskene«. Dette er da et talende Bevis for, at
religiøse Stemningsfantasier uden Selverkjendelse, uden
dybere Forstaaelse af Villieslivet og dets Opgaver, ere
saa langt fra at opklare Selvets Væsen, at de vildsom
krydsende Lysglimt tvertimod forstærke det aandelige
Mørke.

Det er især løsagtige Forestillinger om Guds Naade,
om Syndstilgivelsen og den evige Salighed, der forvirre
de Svage. »Jeg farer nu«, skriver en Selvmorder (Folkebl.
S. 74) »hisset til evig Hvile, hvor jeg, dersom jeg
finder Guds Naade, efter at have bedt om Forladelse
eller Tilgivelse, vil komme til at glæde og fryde mig hos
et høiere altstyrende Væsen der, hvor ingen Synd og
ingen Uret finder Sted\dots{} Jeg vil da skue ned fra det
Høie til Eder Alle her paa Jorden«\dots{} Der er ingen
Tvivl om, at denne noget blødagtige, men, som det synes,
tillige meget barnlige Selvmorder virkelig har havt religiøse
Stemninger; han har vist ogsaa, naar det kom
derpaa an, »kunnet sine Skriftsteder«. Men Et er at
kunne Skriftsteder, et Andet at indøve religiøse Livstanker.
Det kommer, hvad Skriftstederne angaaer, ikke
an paa Mængden, men paa Tilegnelsen. Et eneste Sted,
% --- p. 706 ---
\setcounter{footnote}{0}%
\opage{706}--- det findes i Pauli Brev til Romerne Kap. 14, V. 7-8 ---
er, naar det opfattes i levende Tro og tilegnes i Aand
og Sandhed, et tilforladeligt Skjold og sikkert Værn mod
Fristelse til Selvmord.

Dersom en livlig Syslen med bibelske Forestillinger
og »gudelige Betragtninger« var et paalideligt Mærke
paa sand Religiøsitet, vilde det endda ikke være saa
vanskeligt at gaae i Kamp mod Sindsformørkelsen og
faae Selvmordernes Antal nogenledes indskrænket; men
den religiøse Livsopgave er i allerstrengeste Forstand
en Selvopgave\footnote[1]{De, som i Anledning af de nyeste humane Rørelser i Abstraktionsmetafysiken muligvis antage, at her virkelig er Noget paa Færde, at Tidens Fylde er kommen, da Gud maa give Afkald paa at have en Villie, og at Mennesket, der naturligvis fra sin Side frivillig opgiver Sjælens Udødelighed, kun har, i uendelig Resignation, at forsone sig selv med det Abstrakte, ville, før de altfor troskyldig fordybe sig i den fra Spinoza, Strauss og Feuerbach hentede, ved gjentagen Opvarmning fortyndede og forsaavidt prisbillige Kraftsuppe, --- for en Forsigtigheds Skyld --- gjøre vel i at tage Spørgsmaalet om Selvets Evighed for sig endnu engang, og f. Ex. undersøge om Kritiken, der ellers saa drivende gjendriver Alt, virkelig har „gjendrevet“, hvad S. Kierkegaard i sit Skrift „Sygdommen til Døden“ har oplyst om det syge Selv.}. De, der ikke vide hvad Selvet er,
vide heller ikke hvad Tro er. At blive opmærksom
paa Selvet, dets dunkle Anelser, dets hemmelige Kvaler,
dets Sygdomskriser og Helbredelse, er Opgaven.
Men dette er ogsaa den vanskeligste af alle Opgaver;
thi der er Intet i Verden, som det naturlige Menneske,
just paa Grund af det Selviske, hellere undflyer, end
netop sit eget Selv.

\signature{R. \emph{Nielsen}.}

\end{document}
